\chapter{Frailty}
\index{Frailty}
\label{sec:Frailty}

\section{Overview}


\begin{itemize}[noitemsep]
  \item Frailty is a geriatric syndrome of increased vulnerability to stressors resulting from age-related decline in physiological reserve across multiple organ systems. It affects a significant proportion of older adults. Two major conceptual frameworks guide frailty assessment. The \textbf{Fried Frailty Phenotype} (2001) defines frailty by the presence of three or more of five criteria: unintentional weight loss (>10 lb in one year), self-reported exhaustion, weakness (low grip strength), slow gait speed, and low physical activity
  \item Pre-frailty is defined as one or two criteria. The \textbf{Frailty Index} (Rockwood and Mitnitski) accumulates deficits across multiple domains as a ratio of deficits present to total deficits assessed (typically 30--70 items).
\end{itemize}
\section{Causes}
This condition happens because of age-related decline in physiologic reserve, with sarcopenia as a core biological substrate.     
\section{Risk Factors}

\begin{itemize}[noitemsep]
  \item Genetic predisposition and family history
  \item Advancing age
  \item Lifestyle factors (diet, exercise, smoking, alcohol)
  \item Environmental exposures
  \item Underlying medical conditions
  \item Medication use
\end{itemize}
\section{Signs and Symptoms}

\subsection{Common symptoms}
\begin{itemize}[noitemsep]
  \item You may experience unintentional weight loss, exhaustion, weakness, slow gait speed, and low physical activity. 
  \end{itemize}

\subsection{Emergency warning signs}
\begin{itemize}[noitemsep]
  \item Severe or worsening symptoms of frailty require immediate medical evaluation
\end{itemize}
\section{Progression and Stages}
The condition may progress through different stages of severity over time. Early stages often involve mild or subtle symptoms that may be overlooked. Moderate stages involve more noticeable symptoms affecting daily function. Advanced stages may involve significant impairment and complications.

\begin{itemize}[noitemsep]
  \item The outlook varies from person to person
  \item Frailty is a dynamic process that can improve or worsen over time.
\end{itemize}
\section{Complications}
\begin{itemize}[noitemsep]
  \item Disease progression leading to worsening symptoms
  \item Secondary infections or injuries
  \item Side effects of treatment
  \item Reduced quality of life
  \item Emotional and psychological impact
\end{itemize}
\section{Diagnosis}
Doctors use the Fried Frailty Phenotype (three or more of five criteria) or the Frailty Index (>0.25 indicates frailty). Comprehensive medical history and physical examination Laboratory tests to assess organ function and rule out other conditions Imaging studies to visualize affected structures Specialized testing depending on the affected system Consultation with relevant specialists Monitoring of disease progression over time

\begin{itemize}[noitemsep]
  \item Comprehensive medical history and physical examination
  \item Laboratory tests to assess organ function
  \item Imaging studies to visualize affected structures
  \item Specialized testing depending on the affected system
  \item Consultation with relevant specialists
\end{itemize}
\section{Prevention}
\begin{itemize}[noitemsep]
  \item Maintain a healthy lifestyle with balanced diet and exercise
  \item Avoid known risk factors
  \item Attend regular medical check-ups for early detection
  \item Follow recommended screening guidelines
\end{itemize}
\section{Treatment Options}
Treatment options include resistance and aerobic exercise, protein-calorie supplementation, vitamin D, reduction of polypharmacy, and comprehensive geriatric assessment (CGA). Medications to manage symptoms and address underlying mechanisms Lifestyle modifications and self-care measures Physical therapy and rehabilitation Surgical intervention when indicated Regular monitoring and follow-up care Patient education and support services

\begin{itemize}[noitemsep]
  \item Medications to manage symptoms and address underlying mechanisms
  \item Lifestyle modifications and self-care measures
  \item Physical therapy and rehabilitation
  \item Surgical intervention when indicated
  \item Regular monitoring and follow-up care
\end{itemize}
\section{Living with It}
\begin{itemize}[noitemsep]
  \item Follow your treatment plan as prescribed
  \item Maintain regular follow-up with your healthcare provider
  \item Make necessary lifestyle adjustments
  \item Seek support from family, friends, or support groups
  \item Stay informed about your condition
\end{itemize}
\section{Outlook}
The outlook varies depending on the specific condition, severity at diagnosis, and response to treatment. Early intervention generally leads to better outcomes. Many conditions can be effectively managed with appropriate treatment. Regular follow-up care is important for monitoring and treatment adjustment.
